\chapter{Technische Grundlagen}
\abschnittsautor{M. Musterfrau}

Hier steht alles, was zum Verständnis der praktischen Umsetzung nötig ist --
und nur das. Jede Technologie, die ihr später verwendet, wird hier einmal
erklärt und belegt.

\section{Teilbereich A}

Beim ersten Auftreten schreibt ihr eine Abkürzung aus: \ac{API}. Danach
genügt die Kurzform: \ac{API}. Alle Abkürzungen sammelt ihr in
\texttt{diplomarbeit.tex} im Abkürzungsverzeichnis.

\subsection{Ein Unterabschnitt}

Zitiert wird im IEEE-Stil mit einer Nummer in eckigen
Klammern~\cite{beispiel-artikel}. Mehrere Quellen auf einmal gehen
auch~\cite{beispiel-buch,beispiel-online}.

\subsubsection{Eine Ebene tiefer}

Die vierte Ebene könnt ihr verwenden, aber sparsam. Sie wird nummeriert,
steht aber nicht mehr im Inhaltsverzeichnis. Tiefer gliedert ihr nicht.

\section{Teilbereich B}

Erklärt Technologien in eurer eigenen Formulierung. Wörtlich übernommene
Passagen gehören in Anführungszeichen und mit Seitenangabe belegt, sonst gilt
das als Plagiat.

\pagebreak
\section{Radio Fundemetals}
\subsection{Bandwidth}
\abschnittsautor{J. Tremurici}
LoRa stands for Long Range and this type of wireless communication is designed for sending small data packages over long distances, operating on a battery. It offers sufficient performance for a self governed form of communication and can be extended with extra hops if the network grows.  


\begin{tabular}{p{0.25\linewidth}p{0.25\linewidth}p{0.25\linewidth}p{0.25\linewidth}@{}}
\bfseries Technology & Wireless Communication  & Range & Tx Power\\
\toprule
    Bluetooth & Short range & 10 m & 2.5 mW \\
    Wifi & Short range & 50 m & 80 mW \\
    3G/4G & Cellular & 5 km & 5000 mW \\
    LoRa & LoRa & 
    \begin{itemize}
        \item 2-5 km (urban) 
        \item 5-15 km (rural)
        \item > 15 km (LOS)
    \end{itemize}
    & 500 mW 
\end{tabular}

\pagebreak
\section{Frontend--Backend Communication}
\abschnittsautor{D. E. Leuska}

The dashboard is the frontend of the system. It runs in a browser and presents
sensor readings to the user. The backend runs on the local server and receives
readings from the gateway. PHP can process server requests and return either
complete HTML pages or data for a separate frontend
\cite{php-syntax-html,gfg-php-introduction}. In this design, HTML and CSS
define the page while JavaScript updates its displayed values.

The frontend communicates with the PHP backend by sending HTTP requests to an
application programming interface (API). The API defines which readings can
be requested and how the response is structured. The dashboard can request
recent readings without accessing the database directly. The server can
receive submitted values through PHP request variables, but it must validate
them before use \cite{php-tutorial-forms,w3schools-php-forms}.

JSON is a suitable format for sensor data because it represents structured
records that JavaScript can read. PHP provides \texttt{json\_encode} to
convert server data into JSON \cite{php-json-encode}. Each response should
identify the sensor, measurement time and units. This consistent format lets
the dashboard display different sensors in the same way.

The backend must handle invalid requests and unavailable data explicitly.
Database values should be passed through prepared statements rather than
inserted into SQL text \cite{php-pdo-prepared}. The dashboard should show
when the server cannot be reached or a reading is out of date. Displaying the
measurement time prevents old data from appearing current.
